%%%=============================================================================
%%%  MANUAL DE USO --- Tracker UAI
%%%  Analisis de video y graficos para laboratorios de fisica
%%%  Universidad Adolfo Ibanez / FACH --- Prof. David Aguayo Vera
%%%
%%%  Requiere  uaiestilo.sty  en la misma carpeta.
%%%  Compilar con pdfLaTeX (dos pasadas para indice y referencias).
%%%=============================================================================
\documentclass[11pt,a4paper]{article}

%%-----------------------------------------------------------------------------
%% 1. ESTILO INSTITUCIONAL
%%-----------------------------------------------------------------------------
\usepackage{uaiestilo}

%%-----------------------------------------------------------------------------
%% 2. METADATOS
%%-----------------------------------------------------------------------------
\curso{Universidad Adolfo Ibáñez --- Facultad de Ingeniería y Ciencias}
\cursocorto{Tracker UAI --- Manual de uso}
\title{Tracker UAI}
\subtitle{Manual de uso: análisis de video y gráficos para laboratorios de física}
\author{Prof.\ David Aguayo Vera}
\date{Segundo semestre de 2026}

\uailogos{%
  \uailogo[1.1cm]{logo-uai.png}\hfill%
  \uailogo[1.1cm]{logo-fach.png}}

%%=============================================================================
\begin{document}
%%=============================================================================

\maketitle

\begin{obs}[Cómo usar este manual]
  Este documento explica cómo instalar y usar \textbf{Tracker UAI} en Windows,
  macOS y Linux. Las cajas de color señalan procedimientos paso a paso
  (\textcolor{uaiNaranja}{recetas}), advertencias (\textcolor{uaiRojo}{ojo} y
  errores comunes) e ideas clave. No se requieren conocimientos de programación.
\end{obs}

\tableofcontents
\vspace{4mm}

%%=============================================================================
\section{¿Qué es Tracker UAI?}
%%=============================================================================

Tracker UAI es una aplicación para analizar los videos que los estudiantes graban
en el laboratorio y para construir gráficos de datos experimentales. Permite medir
la posición de un objeto cuadro a cuadro, obtener su velocidad y aceleración,
ajustar modelos físicos y exportar tablas y figuras listas para el informe.

\begin{idea}
  Todo el análisis ocurre en tu computador y en tu navegador. No necesitas cuenta,
  ni internet (salvo la primera instalación), ni instalar Excel, Python o MATLAB
  para usarlo.
\end{idea}

El trabajo se organiza en seis etapas secuenciales para el análisis de video, más
una pestaña independiente de gráficos:

\begin{figure}[htbp]
  \centering
  \begin{tikzpicture}[>=Stealth, node distance=3.5mm,
      paso/.style={draw=uaiBorde, fill=uaiNiebla, rounded corners=2pt,
                   font=\footnotesize\bfseries, text=uaiTinta,
                   minimum height=8mm, inner xsep=5pt}]
    \node[paso] (v)  {1.\,Video};
    \node[paso, right=of v] (c) {2.\,Calibración};
    \node[paso, right=of c] (t) {3.\,Tracking};
    \node[paso, right=of t] (k) {4.\,Cinemática};
    \node[paso, below=6mm of c] (a) {5.\,Ajuste};
    \node[paso, right=of a] (e) {6.\,Exportar};
    \node[paso, fill=uaiVerde!10!white, right=8mm of e] (g) {\faChartLine\ Gráficos};
    \draw[uaiAzul, ->, thick] (v) -- (c);
    \draw[uaiAzul, ->, thick] (c) -- (t);
    \draw[uaiAzul, ->, thick] (t) -- (k);
    \draw[uaiAzul, ->, thick] (k.south) |- (a.north);
    \draw[uaiAzul, ->, thick] (a) -- (e);
  \end{tikzpicture}
  \caption{Flujo de trabajo. La pestaña \emph{Gráficos} es independiente y no
           requiere video.}
  \label{fig:flujo}
\end{figure}

%%=============================================================================
\section{Requisitos e instalación}
%%=============================================================================

Lo único que debes instalar tú es \textbf{Python 3.11 o superior}. El resto de las
librerías las instala la aplicación automáticamente la primera vez que la abres.

\subsection{Instalar Python (una sola vez)}

\begin{receta}[Según tu sistema operativo]
  \begin{itemize}
    \item \textbf{Windows:} descarga Python desde
      \texttt{python.org/downloads} y, en el instalador, marca la casilla
      \textbf{«Add Python to PATH»}. (Alternativa: instalar «Python 3.12» desde la
      Microsoft Store.)
    \item \textbf{macOS:} instala desde \texttt{python.org/downloads} o, con
      Homebrew, ejecuta \texttt{brew install python}.
    \item \textbf{Linux (Ubuntu/Debian):} instala los paquetes base de Python.
  \end{itemize}
\end{receta}

\begin{lstlisting}[caption={Linux: instalación de Python (terminal).}]
sudo apt update
sudo apt install python3 python3-venv python3-pip
\end{lstlisting}

Para comprobar la instalación, abre una terminal y ejecuta \texttt{python -{}-version}
(en Windows) o \texttt{python3 -{}-version} (en macOS/Linux); debe mostrar
\texttt{Python 3.11} o superior.

\begin{errcomun}
  En Windows, olvidar marcar \textbf{«Add Python to PATH»} hace que el programa no
  encuentre Python y no abra. Si te ocurre, reinstala Python marcando esa casilla.
\end{errcomun}

%%=============================================================================
\section{Iniciar el programa}
%%=============================================================================

\begin{receta}[Abrir Tracker UAI]
  \begin{itemize}
    \item \textbf{Windows:} doble clic en \texttt{Iniciar\_Tracker\_UAI.bat}.
    \item \textbf{macOS / Linux:} abre una terminal en la carpeta del programa y
      ejecuta el lanzador (ver abajo).
    \item \textbf{Cualquier sistema (manual):} instala las dependencias y arranca
      con \texttt{streamlit} (ver abajo).
  \end{itemize}
\end{receta}

\begin{lstlisting}[caption={macOS / Linux: lanzador.}]
bash iniciar_tracker.sh
\end{lstlisting}

\begin{lstlisting}[caption={Método manual, cualquier sistema.}]
pip install -r requirements.txt
python -m streamlit run app.py
\end{lstlisting}

La aplicación se abre sola en tu navegador (normalmente en
\texttt{http://localhost:8501}).

\begin{obs}
  La \textbf{primera vez} tarda algunos minutos porque descarga e instala las
  librerías; la ventana de la terminal queda abierta mostrando el avance. Las
  siguientes veces abre en segundos. Para cerrar el programa, cierra esa ventana
  (o presiona \texttt{Ctrl+C} en ella).
\end{obs}

%%=============================================================================
\section{Análisis de video, etapa por etapa}
%%=============================================================================

La barra lateral izquierda contiene el selector de etapas y la carga del video.

\subsection{Etapa 1 — Video}

Sube tu video (mp4, mov, avi, mkv, m4v o webm). Usa el reproductor y los controles
cuadro a cuadro para ubicar el tramo de interés y define el \textbf{intervalo de
procesamiento} (cuadro inicial y final).

\begin{idea}
  Acotar el intervalo hace que todo el análisis sea mucho más rápido: el programa
  trabaja solo sobre el tramo que te interesa, no sobre el video completo.
\end{idea}

\subsection{Etapa 2 — Calibración}

Convierte píxeles a metros. Marca dos puntos de una distancia conocida (por
ejemplo, una regla) y escribe esa distancia en metros; luego marca el origen de
coordenadas.

\begin{receta}[Calibrar]
  \begin{enumerate}
    \item Marca el punto 1 y el punto 2 sobre la distancia conocida.
    \item Escribe la distancia real entre ambos, en metros.
    \item Marca el origen de coordenadas.
    \item (Opcional) Activa \emph{«Ejes con orientación libre»} si el eje $x$ no es
      horizontal (por ejemplo, un plano inclinado) y marca la dirección $+x$.
    \item Pulsa \textbf{Aplicar calibración}.
  \end{enumerate}
\end{receta}

\subsection{Etapa 3 — Tracking (seguimiento del objeto)}

Hay dos modos. En el \textbf{manual}, haces clic sobre el objeto en cada cuadro y
el video avanza solo. En el \textbf{automático}, eliges \emph{Plantilla} (arrastras
un recuadro sobre el objeto), \emph{Color} (haces clic para tomar su color) o
\emph{CSRT}; verás la detección en vivo y, al terminar, todas las detecciones
superpuestas.

\begin{ojo}
  Tras la detección automática puedes \textbf{corregir}: navega cuadro por cuadro y
  haz clic para mover un punto mal detectado, o bórralo. También puedes editar la
  tabla de datos directamente.
\end{ojo}

\subsection{Etapa 4 — Cinemática}

Muestra los gráficos de posición, velocidad, aceleración y trayectoria. Activa el
\textbf{suavizado} si los datos tienen ruido, sobre todo para la aceleración.

\subsection{Etapa 5 — Ajuste de modelos}

Elige un \textbf{preset} según tu laboratorio (gravedad en plano inclinado,
velocidad terminal, colisiones, péndulo cónico, péndulo bifilar, MHD) o el modo
\textbf{genérico}. El programa entrega los parámetros con su \textbf{incertidumbre}
y el coeficiente de determinación $R^{2}$.

\subsection{Etapa 6 — Exportar}

Marca con casillas qué gráficos quieres y descarga un archivo comprimido (ZIP) con
las imágenes (PNG) y el archivo CSV con todos los datos crudos y derivados.

%%=============================================================================
\section{Pestaña «Gráficos» (datos manuales)}
%%=============================================================================

Esta herramienta es \textbf{independiente del video}: sirve para graficar datos
escritos a mano. Está pensada para Física I; por ejemplo, graficar la aceleración
en función de $\sin(\theta)$ en el péndulo simple.

\begin{receta}[Construir un gráfico en pocos pasos]
  \begin{enumerate}
    \item Escribe el título y los nombres de los ejes.
    \item Ingresa tus valores de $x$ e $y$ en la tabla (agrega filas con el botón
      \texttt{+}).
    \item El gráfico de dispersión aparece automáticamente.
    \item Elige el ajuste (lineal o cuadrático): verás la curva, la \textbf{ecuación}
      y el $R^{2}$.
    \item Crea más series (con nombre y color) para comparar conjuntos de datos.
    \item Descarga el gráfico en \textbf{PNG} o \textbf{SVG} para tu informe.
  \end{enumerate}
\end{receta}

\begin{ejemplo}[Péndulo simple: $a$ frente a $\sin\theta$]
  Al graficar la aceleración medida contra $\sin\theta$ y aplicar un ajuste lineal,
  la pendiente corresponde a la aceleración de gravedad $g$, y el $R^{2}$ indica
  qué tan bien la recta describe los datos.
\end{ejemplo}

%%=============================================================================
\section{Consejos para grabar los videos}
%%=============================================================================

\begin{estrategia}[Buenas prácticas de grabación]
  \begin{itemize}
    \item Cámara \textbf{perpendicular} al plano del movimiento y \textbf{fija}
      (trípode o apoyada); esta versión no corrige perspectiva.
    \item La regla de calibración debe estar \textbf{en el mismo plano} que el
      movimiento.
    \item Buena iluminación; objeto bien visible y, si usarás seguimiento por color,
      de un \textbf{color distinto al fondo}.
    \item Para movimientos rápidos, graba en \textbf{cámara lenta} e indica el fps
      real de captura en \emph{Opciones avanzadas}.
  \end{itemize}
\end{estrategia}

%%=============================================================================
\section{Solución de problemas}
%%=============================================================================

\begin{table}[htbp]
  \centering
  \caption{Problemas frecuentes y su solución.}
  \label{tab:problemas}
  \begin{tabular}{@{}p{5.2cm} p{8.3cm}@{}}
    \toprule
    \textbf{Problema} & \textbf{Solución} \\
    \midrule
    «No se encontró Python» & Reinstala Python marcando \textbf{Add to PATH}
      (Windows). \\
    El video \texttt{.mov} no carga o se cuelga & Suele ser el códec HEVC.
      Conviértelo a H.264 con \texttt{ffmpeg}. \\
    El video es muy largo o lento & En la Etapa 1, elige un intervalo corto de
      cuadros. \\
    No aparece el método CSRT & Instalaste \texttt{opencv-python} en vez de
      \texttt{opencv-contrib-python}. \\
    Editar cuadro por cuadro va lento & Acota el intervalo; usa tracking automático
      y corrige solo los puntos malos. \\
    \bottomrule
  \end{tabular}
\end{table}

\begin{alerta}
  Para convertir un video con problemas de códec a un formato compatible (H.264),
  usa la herramienta \texttt{ffmpeg} desde la terminal:
\end{alerta}

\begin{lstlisting}[caption={Convertir un video a H.264.}]
ffmpeg -i entrada.mov -c:v libx264 -pix_fmt yuv420p salida.mp4
\end{lstlisting}

%%=============================================================================
\section{Compartir el programa}
%%=============================================================================

El único requisito en el equipo de destino es tener \textbf{Python 3.11+}
instalado; el resto lo instala el lanzador automáticamente.

\begin{receta}[Distribución]
  \begin{itemize}
    \item \textbf{Windows:} ejecuta \texttt{Crear\_paquete\_para\_compartir.bat}
      para generar un ZIP limpio y envíalo. El estudiante sigue
      \texttt{LEEME\_ESTUDIANTES.txt}.
    \item \textbf{macOS / Linux:} comparte la carpeta (sin la subcarpeta del
      entorno virtual) junto con \texttt{iniciar\_tracker.sh}.
  \end{itemize}
\end{receta}

\begin{resumen}
  \begin{itemize}
    \item Tracker UAI funciona en Windows, macOS y Linux; solo requiere Python 3.11+.
    \item El análisis de video sigue seis etapas: Video, Calibración, Tracking,
      Cinemática, Ajuste y Exportar.
    \item La pestaña \emph{Gráficos} permite analizar datos experimentales sin video.
    \item Todos los resultados se exportan como CSV e imágenes para el informe.
  \end{itemize}
\end{resumen}

\end{document}
